% Chapter 1 -- Introduction
\chapter{Introduction}
\label{ch:introduction}
\lhead{Chapter 1. \emph{Introduction}}

\section{Motivation}
Regulation (EU) 2024/1689 \citep{euaiact2024} makes AI systems that assess the creditworthiness of natural persons high-risk. Its obligations for such stand-alone high-risk systems were to apply from 2~August 2026; Regulation (EU) 2026/1744, the Digital Omnibus on AI, in force since 27~July 2026, postpones them to 2~December 2027 \citep{omnibus}. The stated reason for the postponement is that the standards and support measures providers need in order to demonstrate conformity were not in place --- which is to say that the problem this thesis addresses is the one the legislator has just acknowledged. A bank deploying a foundation model in a lending workflow still needs, by that date, evidence about the model's behaviour that a supervisor will accept. The evidence available today comes in two kinds that do not meet. Benchmarks such as HELM \citep{helm} measure behaviour at scale but are organised by task, so a fairness score names no article, no interpretation and no tolerance. Audit frameworks such as SMACTR \citep{smactr} and the three-layered model of M\"okander et al.\ \citep{mokander} are organised by obligation but stop at process, asking whether documentation exists rather than what the system does when asked to decide a loan.

The preceding case study closed that gap for one domain and, in doing so, exposed a second one. Building the pipeline was tractable; knowing which of its verdicts to believe was not. Quantitative obligations were scored by deterministic arithmetic and are as trustworthy as the arithmetic. Qualitative obligations --- whether an explanation is adequate under Article~13(1) --- were scored by a language model whose agreement with human auditors was never established, because four human annotators following the frozen guideline agreed at chance (Fleiss $\kappa = 0.07$). And when the audit budget was spent adaptively rather than proportionally, coverage rose by up to 17~percentage points and cost halved, while the share of verdicts that survived contact with the full evidence fell from 100\% to 78\%.

These are not three separate shortfalls. They are three views of one boundary: the line past which a machine's verdict can no longer stand in for a human's. Nobody has drawn that line for regulatory auditing, and without it an automated audit is either trusted too much or not used at all.

Nothing in that boundary is specific to credit. A medical-triage model audited against a different instrument raises the same question in the same form: which of its obligations can be scored by machine, at what certified error, and what is lost by stopping early. What credit supplies is a domain where the question can be measured, because the preceding case study already built the corpus, the probes and the 61{,}259 logged responses. The thesis is therefore organised as a procedure plus one worked instance, and is explicit throughout about which of its results are the procedure and which are the instance.

\section{Challenges and Problem Statement}
Three specific problems stand between the existing state of the art and an automated audit whose verdicts can be relied upon.

\paragraph{P1 --- Automated verdicts carry no measured reliability.}\hspace{0.6em}COMPL-AI \citep{complai}, the closest technical bridge between the AI Act and executable benchmarks, maps the Act to existing suites manually and never validates its scoring against human judgement. LLM-as-a-judge research documents position, verbosity and self-preference bias, and JudgeBench \citep{judgebench} reports that the strongest judge reaches only about 64\% on hard objective-correctness pairs. Judges are nevertheless validated once, offline, and never per audit dimension. \emph{Consequence:} a user of an automated audit cannot tell which verdicts to trust.

\paragraph{P2 --- The human reference is assumed rather than measured.}\hspace{0.6em}Validating a judge presupposes a reliable human standard. The case study's pilot shows that this cannot be assumed: raw agreement of 60--80\% coexisted with chance-level $\kappa$, because the label distribution was skewed and the criterion admitted three defensible readings. A judge measured against an unestablished ceiling produces a number that means nothing. \emph{Consequence:} reported judge--human agreement may be uninterpretable, in either direction.

\paragraph{P3 --- Sequential allocation invalidates the verdicts it accelerates.}\hspace{0.6em}Static allocation spends identical budget on settled and unsettled requirements. Adaptive allocation fixes that and introduces optional stopping, which is the classical mechanism by which a sequential procedure manufactures determinations it has not earned --- as the case study's replay demonstrates. No audit framework treats budget allocation as a sequential-inference problem with a spending correction. \emph{Consequence:} an audit is either wasteful, or fast and unsound.

\section{Contributions and Research Questions}
\paragraph{Main research question.}\hspace{0.6em}\emph{Where, across Responsible-AI dimensions, can an automated verdict replace a human one in a clause-traced regulatory audit, at what certified error rate, and at what cost in verdict validity?}

It decomposes into three answerable sub-questions, each with a testable hypothesis. The confirmatory dimension for RQ1 and RQ2 is fixed in advance as transparency (adequacy of explanation under Article~13(1)), the one dimension the case study scored by judge; any further judge-scored dimension is secondary and is reported descriptively, without a confirmatory claim.

\begin{itemize}
\item \textbf{RQ1 (reliability).} What is the human--human agreement ceiling for the transparency dimension, and how closely does a gated language-model judge approach it relative to trivial alternatives? \textbf{H1:} after a decidability-oriented guideline revision, the human ceiling clears $\kappa \geq 0.61$ (lower bootstrap bound) on the primary dimension; the judge, treated as a third rater, is non-inferior to a human rater at a margin of $0.10$; and the judge is superior to a rule-based scorer built directly from the checklist's adequacy conditions.
\item \textbf{RQ2 (certification).} Given a measured ceiling, what selective error can a calibrated gate guarantee, at what coverage, per dimension? \textbf{H2:} on the primary dimension the gate attains selective error $\leq 5\%$ at coverage $\geq 50\%$, with a distribution-free bound.
\item \textbf{RQ3 (sequential cost).} Can an adaptive allocation policy with a sequentially valid stopping rule retain its cost saving while restoring verdict fidelity to the level of a pre-registered static allocation? \textbf{H3:} the corrected policy achieves $\geq 90\%$ of the static policy's verdict fidelity while spending at most 80\% of the static budget.
\end{itemize}

\paragraph{Contributions.}
\begin{itemize}
\item[\textbf{C1}] A domain-independent procedure for establishing the automation boundary: ceiling before judge, judge against ceiling, certified gate, corrected sequential spending --- stated so that it can be executed on any dimension, instrument and domain, with the decision rules and statistics each step requires.
\item[\textbf{C1b}] The filled-in reliability map produced by running that procedure on consumer credit: per dimension, the measured human ceiling, the judge's agreement against it, and the resulting admissibility verdict.
\item[\textbf{C2}] A certified selective-automation gate wired to the judge path, reporting risk--coverage curves with distribution-free bounds and the calibration-set size each guarantee requires.
\item[\textbf{C3}] A sequentially valid audit-allocation procedure, with an empirical comparison against both static and uncorrected-adaptive allocation on cost, coverage and verdict fidelity.
\item[\textbf{C4}] Extended empirical evidence on clause-traced behavioural auditing: a second protected axis and a commercially served model, establishing whether the case study's fairness finding is axis- and family-specific.
\item[\textbf{C5}] An open-source reference implementation and reproducible, clause-traced audit reports, extending the existing repository.
\end{itemize}

\noindent\emph{Conditional contributions} --- delivered if the annotation study does not overrun (Section~\ref{sec:scope}); the thesis is complete without them and is to be assessed on C1--C5.
\begin{itemize}
\item[\textbf{C4b}] A mechanism probe for the observed direction of the fairness effect: whether it tracks the title token itself, a name-only cue, or an explicit statement of gender. This converts a replicated effect into an investigated one and bounds what the title-swap construct measures.
\item[\textbf{C6}] Evidence that the grounding stage is corpus-independent: a second standards corpus parsed and reduced to a signed draft checklist with no code change, reported as unit counts and one signature --- the generality claim made with evidence rather than asserted.
\end{itemize}
